% Nouveau EGA V (cote 25): archive cote 25, site page 140 (PDF page 141).
% typescript. Draft compared: none.
% AI-transcribed by Claude Opus 5.5 (claude-opus-5-5, adaptive thinking, effort high), Message Batches API, batch pilot (msgbatch_0189Nhf1xB615YWnztubtLfc), 2026-11-12; not yet checked by eye.
\scanpage{25/140}
\var{\margin{\uline{nouveau EGA} \uuline{V}}}{\placedmargin{\uline{nouveau EGA} \uuline{V}}}{marginal note placed}{ega25_25_140-v0001}

IV 16

Un nouveau n$^{\text{o}}$ : \uline{Ensemble singulier et supersingulier d'une fonction},
\uline{et critères différentiels}.

Ce n$^{\text{o}}$ sera utilisé au par.20 des sections hyperplanes, mais sa
place naturelle me semble bien au par. 16.

\uline{Définition} 1 Soient X un préschéma régulier, $\hand{\phi}$ une section de $\underline{O}_X$ . Un
point $x \hand{\in} X$ est dit \uline{zéro} \var{\xout{\illegible{3}}}{}{typewriter strike-out}{ega25_25_140-v0002} \uline{singulier} pour $\hand{\phi}$
si on a $\hand{\phi}_x \hand{\in} \underline{m}_x^2$, zéro \uline{supersingulier} si c'est un zéro singulier
et si de plus l'élément de $\underline{m}_x^2/\underline{m}_x^3 \hand{\simeq} \mathrm{Sym}(\underline{m}_x/\underline{m}_x^2)$ qu'il définit,
interprété comme une forme quadratique sur le dual $\underline{t}_x$ de $\underline{m}_x/\underline{m}_x^2$ sur k(x),
est une forme \var{\deleted{non}}{}{deletion}{ega25_25_140-v0003} dégénérée.( Un zéro singulier qui n'est pas supersingu-%
lier est parfois appelé \uline{zéro singulier ordinaire}).

\uline{Remarques} 2 Si $x \hand{\in} V(\hand{\phi})$\hand{,} alors x est un zéro \var{\inserted[typed]{non}}{non}{insertion}{ega25_25_140-v0004} singulier de $\hand{\phi}$ sss
$\hand{\phi}_x \hand{\neq} 0$ et x est un point \var{\xout{\unsure{régxxx}}}{}{typewriter strike-out}{ega25_25_140-v0005} non singulier i.e. régulier de $V(\hand{\phi})$,
i.e. sss x est un point régulier de $V(\hand{\phi})$ et $V(\hand{\phi}) \hand{\neq} X$ au voisinage de
x . \var{\deleted{On peut prouver que si x est singulier pour $\hand{\phi}$ , il est singulier
ordinaire sss $\hand{\phi}_x \hand{\neq} 0$ et la multiplicité de $\underline{O}_{V(\hand{\phi}),x}$ est égale à 2
(ou, comme on dit encore parfois, x est un \uline{point singulier ordinaire}
\uline{du préschéma} $V(\hand{\phi})$ ); donc un $x \hand{\in} V(\hand{\phi})$ est non supersingulier \var{\inserted[typed]{pour $\hand{\phi}$}}{pour $\hand{\phi}$}{insertion}{ega25_25_140-v0007} sss
$\hand{\phi}_x \hand{\neq} 0$ et la multiplicité de $O_{V(\hand{\phi}),x}$ est $\hand{\leq} 2$.}}{}{deletion}{ega25_25_140-v0006}

\uline{Définition} 3. Soient X un préschéma lisse sur un corps k, $\hand{\phi}$ une
section de $\underline{O}_X$ , $x \hand{\in} V(\hand{\phi})$ . On dit que x est un zéro \uline{géométriquement}
\uline{singulier} \var{\xout{\unsure{dexx}}}{}{typewriter strike-out}{ega25_25_140-v0008} (resp. \uline{géométriquement \unsure{supersingulier}} ) de $\hand{\phi}$ rel à k,
si pour toute extension k' de k et tout point \hand{$\xi$} de X à valeurs dans
k\hand{,} localisé en x, la \var{\xout{\illegible{1}}}{}{typewriter strike-out}{ega25_25_140-v0009} point correspondant x' de $X_{k'}$, est un zéro
\var{\xout{\unsure{géom}}}{}{typewriter strike-out}{ega25_25_140-v0010} singulier (resp. supersingulier) de $\hand{\phi}_{k'}$ .

\mspage{\hand{140}}
